\section{混合范式与 Benchmark}
\label{sec:hybrid-benchmark}

\subsection{混合训练三阶段}

\begin{enumerate}[leftmargin=*]
  \item \textbf{阶段 A}（\pathtwo{}）：随机测量 / shadow 自监督预训练（GPT、contrastive、autoencoder）
  \item \textbf{阶段 B}（\pathone{}）：稀疏 labeled oracle 微调
  \item \textbf{阶段 C}（部署）：纯经典 inference 或 surrogate
\end{enumerate}

代表：ShadowGPT~\cite{ShadowGPT2024}、Transformer NQS + Rydberg 实验预训练~\cite{TransformerNQSPretrain2025}、Predictive Surrogates~\cite{PredictiveSurrogates2026}。

\subsection{Benchmark 共识}

\begin{table}[htbp]
  \centering
  \caption{关键 benchmark 结论}
  \small
  \begin{tabular}{@{}llp{6.5cm}@{}}
    \toprule
    文献 & 规模 & 结论 \\
    \midrule
    Bowles 2024 & 160 数据集 & 经典整体优于量子；去纠缠常不降性能 \\
    Schnabel 2025 & 20,000+ 模型 & QKM 高度依赖超参；无 universal winner \\
    Fellner 2026 & 27 时序任务 & VQA 常不如简单经典模型 \\
    QSVM 2026 & 970 实验 & 严格 nested CV 下量子核不 consistently 更优 \\
    \bottomrule
  \end{tabular}
\end{table}

\subsection{项目验收检查清单}

任何 ``quantum data $\rightarrow$ ML'' 项目应报告：
\begin{enumerate}[leftmargin=*]
  \item 经典强基线（同参数量或同训练预算）
  \item Nested CV 或独立 test set
  \item 总 quantum shot budget 与 wall-clock 成本
  \item sim vs hardware 分离结果
  \item 噪声 ablation
  \item measurement protocol + raw outcomes 存档
  \item 负结果 / 无 advantage 情形
\end{enumerate}
